\centerline{\bf \Huge Геометрия масс}

\begin{flushright}
        \scriptsize{}
\end{flushright}

\textit{Материальной точкой} называется пара $(A, m)$, где $A$ --- точка на плоскости, а $m$ --- некоторое число. Число $m$ называется массой этой материальной точки (заметьте, что $m$ не обязательно положительно).

\textit{Центром масс} системы материальных точек $(A_1, m_1), (A_2, m_2), \ldots, (A_n, m_n)$ называется такая точка $Z$, что $m_1 \overrightarrow{ZA_1} + m_2 \overrightarrow{ZA_2} + \ldots + m_n \overrightarrow{ZA_n} = \overrightarrow{0}$.

\problem{0.1}
\textbf{Основная теорема.} Пусть система материальных точек $(A_1, m_1), (A_2, m_2), \ldots, (A_n, m_n)$ такова, что $m_1 + m_2 + \ldots + m_n \neq 0$. Тогда центр масс этой системы существует и единственен, причем для любой точки $O$ верно следующее равенство:
\[
\overrightarrow{OZ} = \frac{m_1 \overrightarrow{OA_1} + m_2 \overrightarrow{OA_2} + \dots + m_n \overrightarrow{OA_n}}{m_1 + m_2 + \ldots + m_n}
\]

\problem{0.2}
\textbf{Правило рычага.} Центр масс $Z$ двух материальных точек $(A, m_1)$ и $(B, m_2)$ с неотрицательными массами расположен на отрезке $AB$, причем $BZ : AZ = m_1 : m_2$. Что произойдет, если убрать условие неотрицательности масс?

\problem{0.3} 
\textbf{Правило группировки.} Пусть $O$ --- центр масс первых $k$ материальных точек системы $(A_1, m_1), (A_2, m_2), \ldots, (A_n, m_n)$. Тогда центр масс всей системы совпадает с центром масс системы $(O, m_1 + \ldots + m_k), (A_{k+1}, m_{k+1}), \ldots, (A_n, m_n)$.

\problem{1}
Пусть точка $X$ лежит внутри треугольника $ABC.$ Докажите, что если на вершины $A, B, C$ повесить массы $S_{BXC}, S_{AXC}, S_{AXB}$ соответственно, то точка $X$ будет центром масс этой системы.

\problem{2}
Дан треугольник со сторонами $a, b, c$ и углами $\alpha, \beta, \gamma$ Какие массы надо поместить в вершины этого треугольника, чтобы центр масс полученной системы материальных точек оказался (а) в точке пересечения медиан; (б) в точке пересечения биссектрис; (в) в точке пересечения высот; (г) в центре описанной окружности; (д) в точке Жергона; (e) в точке Нагеля?

\newpage

\hspace{3mm}

\problem{3}
На сторонах $AB$ и $AC$ треугольника $ABC$ выбраны точки $D$ и $E$ соответственно. $X$ $-$ точка пересечения отрезков $BE$ и $CD$. В точке $A$ находится масса $6$. Какие массы надо поместить в точки $B$ и $C$, чтобы центр масс треугольника попал в точку $X$, если $AD : DB = 1 : 1, AE : EC = 3 : 1?$

\problem{4}
На сторонах $AB$ и $BC$ треугольника $ABC$ расположены точки $M$ и $N$ соответственно, причём $AM : MB = 3 : 5, BN : NC = 1 : 4.$ Прямые $CM$ и $AN$ пересекаются в точке $O$. Найдите отношения $AO : ON$ и $CO : OM.$

\problem{5}
На сторонах $AB, BC$ и $AC$ треугольника $ABC$ взяты соответственно точки $M, N$ и $K$ так, что $AM : MB = 2 : 3, AK : KC = 2 : 1, BN : NC = 1 : 2.$ В каком отношении прямая $MK$ делит отрезок $AN$?

\problem{6}
Около окружности описан четырёхугольник $ABCD$, касающийся окружности в точках $M, N, P, Q.$ Известно, что длины отрезков касательных, проведённых из вершин $A, B, C, D$ к окружности равны $a, b, c, d$ соответственно. В каком отношении делится каждый из отрезков $MP$ и $NQ$ точкой их пересечения?

\problem{7}
\textbf{Теорема Чевы.}
Определим чевиану как отрезок, соединяющий вершину треугольника с некоторой точкой на противоположной стороне.

Три чевианы ${\displaystyle AA',BB',CC'}$ треугольника ${\displaystyle ABC}$ проходят через одну точку тогда и только тогда, когда: 
\[{\displaystyle {\frac {BA'\cdot CB'\cdot AC'}{A'C\cdot B'A\cdot C'B}}=1}\]

\problem{8}
\textbf{Теорема Менелая.}
Если точки $A', B'$ и $C'$ лежат соответственно на сторонах $BC, CA$ и $AB$ треугольника $\triangle ABC$ или на их продолжениях, то они коллинеарны тогда и только тогда, когда
\[
\frac{AB'}{B'C'} \cdot \frac{CA'}{A'B} \cdot \frac{BC'}{C'A} = -1.
\]
где $\frac{AB'}{B'C'}, \frac{CA'}{A'B}$ и $\frac{BC'}{C'A}$ обозначают отношения направленных отрезков.

\problem{9}
\textbf{Теорема Ван-Обеля.}
Если прямые $AP, BP, CP$ пересекают соответственно прямые $BC, CA$ и $AB$, содержащие стороны треугольника $ABC$ соответственно в точках $A_1, B_1$ и $C_1$, то имеет место равенство отношений направленных отрезков:
\[
\frac{AC_1}{C_1B} + \frac{AB_1}{B_1C} = \frac{AP}{PA_1}.
\]